% GENERATED FILE. Edit canonical lessons, then run npm run generate:books.
% canonical-source-hash: fnv1a64:29e9f38dff50faa4
% canonical-lessons: RU-C365-gosudarstvo, RU-C365-pravitelstvo, RU-C365-vlast, RU-C365-politika, RU-C365-politik

\chapter{The state and power}
\label{ch:ru-the-state-and-power}
\hlchaptermodality{365}

\begin{chapteropening}
\textbf{By the end of this chapter:} \emph{I can talk about the state, the government, power and politics, and about a politician.}

Complete the last lesson of chapter 365: I can talk about the state, the government, power and politics, and about a politician.
\end{chapteropening}

\section[\texorpdfstring{\ru{государство} (gosudárstvo)}{государство (gosudárstvo)}]{\ru{государство} (gosudárstvo) — the state, a country}

\label{lesson:RU-C365-gosudarstvo}



\begin{quote}
Before the new one: say the Russian for social, then the Russian for an inhabitant.
\end{quote}

\subsection*{You'll want to know: \ru{государство}}
\textbf{\ru{государство}} (\emph{gosudárstvo}) — \textquotedblleft{}the state, a country\textquotedblright{}.

\begin{quote}
\ru{Государство} \ru{помогает} \ru{пенсионерам}.
\end{quote}

\begin{quote}
\emph{Gosudárstvo pomogáyet pensionéram.}
\end{quote}

\begin{quote}
The state helps pensioners.
\end{quote}

Neuter. \emph{\ru{Государственный}} — state, public: \emph{\ru{государственный} \ru{язык}}.

\begin{cousinweb}
\textbf{\ru{государство}} is built on \emph{\ru{государь}} \textquotedblleft{}a sovereign\textquotedblright{}, which is usually traced to older \emph{\ru{господарь}} \textquotedblleft{}a lord, a master\textquotedblright{}, from \emph{\ru{господь}} \textquotedblleft{}lord\textquotedblright{}.
\end{cousinweb}

\begin{grammarlens}[title={What you've built}]
One more word for telling, explaining and describing at B1.
\end{grammarlens}

\subsection*{Your turn}
\begin{itemize}
\raggedright
  \item \emph{Say it:} \emph{gosudárstvo}
  \item \emph{Say it:} \emph{gosudárstvo}, once more
  \item \emph{Say it:} a sentence of your own with \emph{gosudárstvo}
  \item \emph{Recall:} say the Russian for social, then the Russian for an inhabitant, then say \emph{gosudárstvo} again
\end{itemize}

\begin{tcolorbox}[breakable,colback=teal!4,colframe=teal!35!black,title={Before you move on}]
What does \ru{государство} mean? (\textquotedblleft{}The state, a country\textquotedblright{}.) Say it once more.
\end{tcolorbox}
\section[\texorpdfstring{\ru{правительство} (pravítelstvo)}{правительство (pravítelstvo)}]{\ru{правительство} (pravítelstvo) — the government}

\label{lesson:RU-C365-pravitelstvo}



\begin{quote}
Before the new one: say the Russian for an inhabitant, then the Russian for the state.
\end{quote}

\subsection*{You'll want to know: \ru{правительство}}
\textbf{\ru{правительство}} (\emph{pravítelstvo}) — \textquotedblleft{}the government\textquotedblright{}.

\begin{quote}
\ru{Правительство} \ru{обещало} \ru{повысить} \ru{пенсии}.
\end{quote}

\begin{quote}
\emph{Pravítelstvo obeshchálo povýsit pénsii.}
\end{quote}

\begin{quote}
The government promised to raise pensions.
\end{quote}

Neuter. The cabinet that runs a country; \emph{\ru{власти}} is broader.

\begin{cousinweb}
\textbf{\ru{правительство}} is built on \emph{\ru{правитель}} \textquotedblleft{}a ruler\textquotedblright{}, from \emph{\ru{править}} \textquotedblleft{}to rule, to steer\textquotedblright{}, from \emph{\ru{правый}} \textquotedblleft{}right\textquotedblright{}, Proto-Slavic \emph{*pravŭ} \textquotedblleft{}right, correct\textquotedblright{} (EtymDB).
\end{cousinweb}

\begin{grammarlens}[title={What you've built}]
One more word for telling, explaining and describing at B1.
\end{grammarlens}

\subsection*{Your turn}
\begin{itemize}
\raggedright
  \item \emph{Say it:} \emph{pravítelstvo}
  \item \emph{Say it:} \emph{pravítelstvo}, once more
  \item \emph{Say it:} a sentence of your own with \emph{pravítelstvo}
  \item \emph{Recall:} say the Russian for an inhabitant, then the Russian for the state, then say \emph{pravítelstvo} again
\end{itemize}

\begin{tcolorbox}[breakable,colback=teal!4,colframe=teal!35!black,title={Before you move on}]
What does \ru{правительство} mean? (\textquotedblleft{}The government\textquotedblright{}.) Say it once more.
\end{tcolorbox}
\section[\texorpdfstring{\ru{власть} (vlast)}{власть (vlast)}]{\ru{власть} (vlast) — power, authority}

\label{lesson:RU-C365-vlast}



\begin{quote}
Before the new one: say the Russian for the state, then the Russian for the government.
\end{quote}

\subsection*{You'll want to know: \ru{власть}}
\textbf{\ru{власть}} (\emph{vlast}) — \textquotedblleft{}power, authority\textquotedblright{}.

\begin{quote}
\ru{В} \ru{этой} \ru{стране} \ru{власть} \ru{часто} \ru{меняется}.
\end{quote}

\begin{quote}
\emph{V étoy strané vlast chásto menyáyetsya.}
\end{quote}

\begin{quote}
In this country power often changes hands.
\end{quote}

Feminine. The plural \emph{\ru{власти}} means the authorities.

\begin{cousinweb}
\textbf{\ru{власть}} \textquotedblleft{}power, authority\textquotedblright{} has the Church Slavonic \emph{-\ru{ла}-}: Church Slavonic \emph{vlastĭ} and native \emph{\ru{волость}} both come from Proto-Slavic \emph{*volstĭ} \textquotedblleft{}rule, power\textquotedblright{} (EtymDB). \emph{\ru{Волость}} was once a district under someone’s rule.
\end{cousinweb}

\begin{grammarlens}[title={What you've built}]
One more word for telling, explaining and describing at B1.
\end{grammarlens}

\subsection*{Your turn}
\begin{itemize}
\raggedright
  \item \emph{Say it:} \emph{vlast}
  \item \emph{Say it:} \emph{vlast}, once more
  \item \emph{Say it:} a sentence of your own with \emph{vlast}
  \item \emph{Recall:} say the Russian for the state, then the Russian for the government, then say \emph{vlast} again
\end{itemize}

\begin{tcolorbox}[breakable,colback=teal!4,colframe=teal!35!black,title={Before you move on}]
What does \ru{власть} mean? (\textquotedblleft{}Power, authority\textquotedblright{}.) Say it once more.
\end{tcolorbox}
\section[\texorpdfstring{\ru{политика} (polítika)}{политика (polítika)}]{\ru{политика} (polítika) — politics; a policy}

\label{lesson:RU-C365-politika}



\begin{quote}
Before the new one: say the Russian for the government, then the Russian for power.
\end{quote}

\subsection*{You'll want to know: \ru{политика}}
\textbf{\ru{политика}} (\emph{polítika}) — \textquotedblleft{}politics; a policy\textquotedblright{}.

\begin{quote}
\ru{Я} \ru{не} \ru{интересуюсь} \ru{политикой}.
\end{quote}

\begin{quote}
\emph{Ya ne interesúyus polítikoy.}
\end{quote}

\begin{quote}
I am not interested in politics.
\end{quote}

Feminine. Both politics and a policy: \emph{\ru{политика} \ru{компании}}.

\begin{cousinweb}
\textbf{\ru{политика}} is usually traced, through western languages, to Greek \emph{politik\'{\={e}}} \textquotedblleft{}the art of running a state\textquotedblright{}, built on \emph{pólis} \textquotedblleft{}a city, a state\textquotedblright{}.
\end{cousinweb}

\begin{grammarlens}[title={What you've built}]
One more word for telling, explaining and describing at B1.
\end{grammarlens}

\subsection*{Your turn}
\begin{itemize}
\raggedright
  \item \emph{Say it:} \emph{polítika}
  \item \emph{Say it:} \emph{polítika}, once more
  \item \emph{Say it:} a sentence of your own with \emph{polítika}
  \item \emph{Recall:} say the Russian for the government, then the Russian for power, then say \emph{polítika} again
\end{itemize}

\begin{tcolorbox}[breakable,colback=teal!4,colframe=teal!35!black,title={Before you move on}]
What does \ru{политика} mean? (\textquotedblleft{}Politics; a policy\textquotedblright{}.) Say it once more.
\end{tcolorbox}
\section[\texorpdfstring{\ru{политик} (polítik)}{политик (polítik)}]{\ru{политик} (polítik) — a politician}

\label{lesson:RU-C365-politik}



\begin{quote}
Before the new one: say the Russian for power, then the Russian for politics.
\end{quote}

\subsection*{You'll want to know: \ru{политик}}
\textbf{\ru{политик}} (\emph{polítik}) — \textquotedblleft{}a politician\textquotedblright{}.

\begin{quote}
\ru{Этот} \ru{политик} \ru{часто} \ru{выступает} \ru{по} \ru{телевизору}.
\end{quote}

\begin{quote}
\emph{Étot polítik chásto vystupáyet po televízoru.}
\end{quote}

\begin{quote}
This politician often speaks on television.
\end{quote}

Masculine, also of a woman. Not to be confused with \emph{\ru{политика}}, politics.

\begin{cousinweb}
\textbf{\ru{политик}} is usually traced, through German, to Greek \emph{politikós} \textquotedblleft{}of the citizens, of the state\textquotedblright{}, built on \emph{pólis} \textquotedblleft{}a city, a state\textquotedblright{}.
\end{cousinweb}

\begin{grammarlens}[title={What you've built}]
One more word for telling, explaining and describing at B1.
\end{grammarlens}

\subsection*{Your turn}
\begin{itemize}
\raggedright
  \item \emph{Say it:} \emph{polítik}
  \item \emph{Say it:} \emph{polítik}, once more
  \item \emph{Say it:} a sentence of your own with \emph{polítik}
  \item \emph{Recall:} say the Russian for power, then the Russian for politics, then say \emph{polítik} again
\end{itemize}

\begin{tcolorbox}[breakable,colback=teal!4,colframe=teal!35!black,title={Before you move on}]
What does \ru{политик} mean? (\textquotedblleft{}A politician\textquotedblright{}.) Say it once more.
\end{tcolorbox}
